% ==============================================================================
% Fichier : front/01_avant_propos.tex
% Titre   : Avant-propos - GCI
% Auteur  : MINKA MI NGUIDJOÏ Thierry Emmanuel
% Version : 1.0 -- octobre 2026
% ==============================================================================

\chapter*{Avant-propos}
\addcontentsline{toc}{chapter}{Avant-propos}
\markboth{Avant-propos}{Avant-propos}

\epigraph{%
    « Spécifier les droits de décision et le cadre de responsabilité qui
    encouragent les comportements souhaitables dans l'usage des technologies
    de l'information. »%
}{Peter \textsc{Weill} et Jeanne \textsc{Ross}, définition de la
gouvernance des SI (2004, notre traduction)}

\vspace{0.5ex}

Ce cours repose sur une double conviction. La première est que la
gouvernance des systèmes d'information n'est plus une affaire de
spécialistes. Au Cameroun comme dans toute l'Afrique centrale, les banques,
les opérateurs, les administrations et les jeunes entreprises du numérique
confient désormais leur activité à des systèmes dont la défaillance arrête
tout : les paiements, les salaires, les ports, les soins. Quand un système
s'arrête, la question qui suit n'est jamais seulement technique. Elle est
toujours la même : \emph{qui avait décidé, et qui devait rendre compte~?}

La seconde conviction est que cette discipline ne s'importe pas telle
quelle. COBIT, ITIL et les normes ISO sont des acquis précieux, mais ils
ont été pensés dans des contextes où les régulateurs, les marchés, les
compétences et les contraintes budgétaires ne sont pas les nôtres. Les
appliquer en Afrique suppose de les relire à la lumière du droit
camerounais et communautaire (CEMAC, OHADA), des attentes de nos
régulateurs et de la réalité de nos organisations.

\bigskip

\textbf{Ce que ce cours n'est pas.} Il n'est ni un résumé des publications
d'ISACA ou de PeopleCert, ni une préparation à une certification. Les
référentiels y sont présentés pour ce qu'ils font bien, cités avec
exactitude, et discutés là où leur transposition pose problème.

\bigskip

\textbf{Ce que ce cours est.} Un traité en six parties, précédées d'un
prologue, qui va du cas concret vers la théorie, de la règle vers la
pratique, et du cas simple vers la recherche. La Partie~I pose les
fondements de la gouvernance des SI. La Partie~II est consacrée à
COBIT~2019, la Partie~III au cadre normatif et juridique, la Partie~IV à
ITIL~4 et aux processus, la Partie~V aux défis de l'agilité, du nuage et de
l'intelligence artificielle, et la Partie~VI à la recherche. Quatre
organisations fictives, une banque, un opérateur portuaire, une
administration et une jeune entreprise de paiement mobile, traversent le
cours et servent de terrain aux exercices.

\bigskip

\textbf{Comment utiliser ce cours.} Chaque chapitre s'ouvre sur une \emph{boîte
objectifs} précisant les compétences attendues, et se ferme sur une \emph{boîte
résumé}. Les exercices sont gradués en trois niveaux : \niveaubase\ fondamental,
\niveaumoyen\ intermédiaire, \niveauexpert\ expert. Chaque chapitre est
accompagné d'une fiche d'auto-évaluation sur NDAP SUKULU, et dix
laboratoires pratiques d'une heure complètent les 30~heures de présentiel
comme travail personnel.

\bigskip

\begin{flushright}
    \textit{\auteurcours}\\
    \textit{Yaoundé, octobre 2026}
\end{flushright}
